\section{模型评价与推广}

\subsection{模型优点}

\textbf{（1）统一框架、一以贯之。} 本文将四个看似递进的问题统一为含储能的单母线能量流线性规划，四问的差别仅在于"进入模型的序列是确定值、预测量还是带更新的预测量"以及"偏差如何计价"。全部模型共享同一套物理约束（\cref{eq:balance}--\cref{eq:bounds}）与同一条守恒关系（\cref{eq:identity}），避免了"一问一模型、彼此割裂"的常见问题。

\textbf{（2）可解性与精确性。} 四问均为线性规划：决策变量连续、目标函数为凸分段线性、约束线性，故可用单纯形/内点法求得\textbf{全局最优解}而无需启发式搜索。本文进一步证明了最优解中充放电不会同时发生（\cref{eq:mutex}），从而说明无须引入 0-1 变量即可保持 LP 结构。问题2--4 的最复杂形式（多阶段随机规划）仍是大规模线性规划，可精确求解。

\textbf{（3）三条理论线索支撑结果解释。} 其一，储量影子价格给出的\textbf{阈值策略}（\cref{eq:threshold}）解释了最优计划的时段结构；其二，报童模型给出的\textbf{覆盖分位点区间}（\cref{eq:p2-newsvendor}）与调整机制下的\textbf{中位数准则}（\cref{eq:p3-median}）解释了"计划应留多少裕度"，并揭示了费率结构决定分位点的规律；其三，\textbf{信息价值分解}（\cref{eq:p3-value}）把"是否需要更密的预报"转化为可定量比较的价值量。三条线索使论文不仅给出数值，还给出"为什么是这样"的解释。

\textbf{（4）结果可复现、可核验。} 全部结果由脚本一键重跑，随机种子固定；核心结论经独立算法（母线侧布局 LP 与动态规划）与独立结算复算交叉验证，$40$ 项检查全部通过。

\subsection{模型不足}

\textbf{（1）终端电量口径。} 本文采用日周期条件 $E_{144}=E_0$，对问题4（电价逐日波动）会略微低估跨日套利的收益，结果偏保守；跨日滚动衔接或全时段联合优化可给出一致性更好的结果，但需要额外的终端价值估计。

\textbf{（2）储能执行采用"开环计划"。} 模型假设储能在日内严格按 0:00 计划执行、缺额以紧急购电填补，未考虑"实时调整充放电以吸收实际光伏波动"的闭环控制。若允许储能实时响应，紧急购电可进一步下降，但这会改变问题2 的费用结构，需与题目给定的执行机制一并讨论。

\textbf{（3）负载侧无预报产品。} 题目只提供光伏预报，负载只能由历史数据外推，其误差（RMSE $955.9$ \si{kW}）是本文费用的主要来源，且在该信息条件下不可改善。

\textbf{（4）其他简化。} 未考虑储能循环寿命与折旧成本、未考虑配电网潮流与电压约束、未区分日前市场与实时市场的双结算机制、未考虑光伏出力与价格的联合随机性（本文以情景/边界方式处理）。这些简化在竞赛给定的数据条件下是必要的，但会影响策略在真实系统中的适用性。

\subsection{模型推广}

\textbf{（1）方法层面的推广。} 本文提出的"统一线性规划内核 $+$ 安全裕度 $+$ 滚动重优化 $+$ 信息价值分解"四件套，可直接迁移到任何"含储能的时序购售电/调度"问题：园区综合能源系统、电动汽车充电站与光储充一体化场站、数据中心后备储能与峰谷套利、用户侧需求响应等，只需替换母线平衡中的电源与负荷构成。

\textbf{（2）问题层面的推广。} 若把单母线扩展为多节点配电网，可在 \cref{eq:balance} 中增加节点功率平衡与线路潮流约束（仍为线性），把问题变为多微网协同优化；若引入需求响应，可把负载 $L_t$ 变为决策变量并加入用电效用函数；若引入储能寿命模型，可在目标中加入与充放电深度相关的折旧项（得到分段线性或二次目标）。

\textbf{（3）机制层面的推广。} 本文的结论对市场机制设计亦有参考价值：由于调整罚金会抵消预报更新的收益（\S5.6），若不希望市场主体通过频繁调整套取信息红利，应保持费率结构的对称性；若希望鼓励主体利用更新信息改善计划，则应降低调整费率或对"因预报修订而被动调整"的部分予以豁免。
